% from ega4-2-fr.tex
\oldpage[IV-2]{214}

A) Pour un anneau \emph{local} noethérien $A$, les propriétés de $A$ se
« transmettent-elles » à son \emph{complété} $\widehat A$ ? Par exemple, si
$A$ est réduit (resp. intègre, resp. intègre et intégralement clos), en est-il
de même de $\widehat A$ ? La plupart de ces questions sont liées aux propriétés
locales des \emph{fibres formelles} de $A$ (rappelons que ce sont les fibres du
morphisme canonique
$\operatorname{Spec}(\widehat A)\to\operatorname{Spec}(A)$). Une exception est
formée par les propriétés liées à la notion de dimension, par exemple la
propriété d'être équidimensionnel; c'est alors la \emph{condition des chaînes}
et ses divers raffinements qui jouent le rôle essentiel.

B) Pour un préschéma localement noethérien $X$ (en particulier pour un schéma
affine $\operatorname{Spec}(A)$), l'ensemble des $x\in X$ où l'anneau local
$\mathcal O_x$ possède une certaine propriété (par exemple être intégralement
clos, ou de Cohen-Macaulay, ou régulier) est-il \emph{ouvert} ?

C) Pour un anneau \emph{intègre} $A$, la fermeture intégrale de $A$ dans une
extension finie de son corps des fractions est-elle un $A$-module \emph{de type
fini} ? Bien entendu, cette question peut se traduire pour les préschémas
noethériens (\textbf{II}, 6.3).

Les problèmes de type B) ou C) peuvent aussi être posés pour des anneaux
\emph{locaux}, mais il ne suffit pas en général qu'ils soient résolus
affirmativement pour tout anneau local $A_{\mathfrak p}$ d'un anneau $A$ pour
qu'ils le soient pour $A$ (cf. (6.13.6)).

Soulignons en outre que dans l'étude de ces problèmes, nous nous sommes
systématiquement préoccupés de savoir si la réponse affirmative à l'un d'eux
est \emph{stable} pour les deux opérations les plus importantes de l'Algèbre
commutative : la localisation et le passage à une algèbre de type fini.

Les résultats obtenus dans l'étude de ces problèmes amènent à dégager la
définition d'une classe d'anneaux noethériens dont le comportement à cet égard
est le meilleur possible :
\end{env}

\begin{definition}[7.8.2]
\EGAReaderTarget{ega4.c04.def.7-8-2}{IV.2.s7.8.2-fr}
On dit qu'un anneau $A$ est excellent s'il est noethérien et s'il
vérifie les conditions suivantes :
\begin{enumerate}
\item[\emph{(i)}] $A$ est universellement caténaire (ou, ce qui revient au même
(5.6.3, (i)), pour tout idéal premier $\mathfrak p$ de $A$,
$A_{\mathfrak p}$ est universellement caténaire).
\item[\emph{(ii)}] Pour tout idéal premier $\mathfrak p$ de $A$, les fibres
formelles de $A_{\mathfrak p}$ sont géométriquement régulières.
\item[\emph{(iii)}] Pour tout anneau quotient intègre $B$ de $A$ et toute
extension finie radicielle $K'$ du corps des fractions $K$ de $B$, il existe
une sous-$B$-algèbre finie $B'$ de $K'$, contenant $B$, ayant $K'$ pour corps
des fractions, et telle que l'ensemble des points réguliers de
$\operatorname{Spec}(B')$ contienne un ouvert non vide.
\end{enumerate}
\upshape
Avec cette terminologie, l'essentiel des résultats du \S~7 et d'une partie du
\S~6 se résume alors comme suit :
\end{definition}

\begin{scholium}[7.8.3]
\EGAReaderTarget{ega4.c04.statement.7-8-3}{IV.2.s7.8.3-fr}
(i) Dans les conditions (i) et (ii) de la définition (7.8.2) on peut se borner
aux idéaux $\mathfrak p$ maximaux dans $A$. Pour qu'un anneau local noethérien
soit excellent, il faut et

